% =====================================================================
\section{Overview and chapter map}
\label{sec:ch7_overview}
% =====================================================================

Chapter~\ref{chap:eda} characterised the daily pattern of ED arrivals at Steve Biko Academic Hospital, and Chapter~\ref{chap:models} produced a daily forecast at a test MAPE of 12.63 percent. Those numbers are useful to the hospital only when they translate into operational decisions, and the decisions sit in two layers: the pharmacy supply chain that holds the medicines tomorrow's patients will consume, and the nurse roster that defines who will be on duty when they arrive. This chapter asks two questions about those two layers.

The first question is whether the chaos that the published South African public-sector literature documents in pharmacy supply and in nurse rostering can be reproduced under a calibrated simulation. If the answer is yes, the chapter has a credible model of the operational environment the forecast must serve. The second question is how much of that chaos a forecast-driven replenishment policy can recover, so the chapter can quantify the operational value of Chapter~\ref{chap:models}.

A structural finding sits between those two questions and is reported in its own right. The same calibrated simulation makes the difference between actual coverage and lawful coverage numerically visible: the roster reaches approximately 97~percent coverage only because filled nurses work a mean of 58 weekly hours, which is 129~percent of the 45-hour weekly cap set by Section~9 of the Basic Conditions of Employment Act \parencite{bcea1997}. Holding the same 23 filled posts to the lawful 45-hour cap leaves the same demand at approximately 75~percent coverage, a shortfall of roughly 7 posts against the lawful establishment. The gap is the operational expression of the staff-shortage pattern documented across South African public hospitals \parencite{malakoane2020public, abrahams2022}, and is one of the chapter's headline findings.

The chapter is organised in four parts. Section~\ref{sec:ch7_methodology} sets out the methodological approach, a calibrated simulation anchored to peer-reviewed South African literature rather than a collected dataset. Sections~\ref{sec:ch7_inventory_model} and~\ref{sec:ch7_scheduling_model} describe the inventory and scheduling models in turn, naming the Python library stack and the literature source behind every stochastic mechanism. Sections~\ref{sec:ch7_inputs} and~\ref{sec:ch7_experimental_design} fix the inputs and the experimental protocol (30 independent replications, 30-day warm-up, common random numbers). Sections~\ref{sec:ch7_calibration} to~\ref{sec:ch7_sensitivity} report the calibration of observed against pre-registered targets, the goodness-of-fit checks, and a one-factor sensitivity sweep. Section~\ref{sec:ch7_baseline_vs_forecast} describes the baseline-versus-forecast comparison that is the chapter's headline experiment and is flagged as documented future work in Section~\ref{sec:ch7_limitations}. Sections~\ref{sec:ch7_discussion} to~\ref{sec:ch7_reproducibility} close the chapter with discussion, limitations, and the reproducibility statement.
